% Part: turing-machines
% Chapter: undecidability
% Section: verification

\documentclass[../../../include/open-logic-section]{subfiles}

\begin{document}

\olfileid{tur}{und}{ver}
\olsection{التحقق من صحة التمثيل}

\begin{explain}
صحة هذا التمثيل تتم بإثبات جهتين. الأولى: إذا توقفت $M$ عند
الدخل~$w$، صدقت $!T(M,w)\lif !E(M,w)$ منطقيًا. والثانية عكسها:
إذا كانت $!T(M,w)\lif !E(M,w)$ صادقة منطقيًا، لزم أن تتوقف
$M$ فعلًا في نهاية المطاف عند الدخل~$w$.

ولكل جهة طريق غير طريق الأخرى. ففي الأولى نريد إثبات الصدق
المنطقي ل!!a{sentence} من الرتبة الأولى، هي
$!T(M,w)\lif !E(M,w)$. وأيسر ذلك أن نقيم لها !!a{derivation}.
غير أن الحكم عام في $M$ و$w$؛ فليس المطلوب !!{sentence}
واحدة نقيم لها !!a{derivation}، بل جمل لا تتناهى. فلا بد من
حجة استقرائية تبيّن أنه، أيًا كانت $M$ و~$w$، وأيًا كان عدد
خطوات $M$ قبل توقفها عند الدخل~$w$، يوجد !!a{derivation}
لـ$!T(M,w)\lif !E(M,w)$.

ومحل الاستقراء عدد الخطوات التي تقطعها $M$ إلى تهيئة التوقف.
فنثبت لكل خطوة~$n$ من تشغيل $M$ عند الدخل~$w$، إلى أول توقف،
أن $!T(M,w)\Entails !C(M,w,n)$؛ حيث $!C(M,w,n)$ وصف صحيح
لتهيئة $M$ عند~$w$ بعد~$n$ خطوات. فإذا توقفت $M$ عند
الدخل~$w$ بعد $n$ خطوة، كانت $!C(M,w,n)$ وصفًا لتهيئة توقف.
ونثبت أيضًا أن $!C(M,w,n)\Entails !E(M,w)$ متى كانت
$!C(M,w,n)$ كذلك. فإذا توقفت $M$ عند الدخل~$w$، وجد~$n$
تكون عنده $M$ في تهيئة توقف. وحينئذ
$!T(M,w)\Entails !C(M,w,n)$، وتصف $!C(M,w,n)$ ذلك التوقف؛
فينضم إلى هذا $!C(M,w,n)\Entails !E(M,w)$، وينتج
$!T(M,w)\Entails !E(M,w)$، أي
$\Entails !T(M,w)\lif !E(M,w)$.

وأما العكس فنبدأ فيه بفرض $\Entails !T(M,w)\lif !E(M,w)$،
ونطلب إثبات توقف $M$ عند الدخل~$w$. والفرض يفيد
$!T(M,w)\Entails !E(M,w)$؛ أي تصدق $!E(M,w)$ في كل
!!{structure} تصدق فيها $!T(M,w)$. فننشئ
!!a{structure}~$\Struct{M}$ تحقق $!T(M,w)$، مجالها $\Nat$؛ ومن تهيئات~$M$ في تشغيلها عند الدخل~$w$
نأخذ تفسيرات $\Obj Q_q$ و$\Obj S_\sigma$. فيكون، مثلًا،
$\Sat{M}{\Obj Q_q(\num{m},\num{n})}$ إذا وفقط إذا كانت $M$
عند الدخل~$w$، بعد $n$ خطوة، في الحالة~$q$ ماسحةً الخانة~$m$.
فلدينا $!T(M,w)\Entails !E(M,w)$ بالفرض،
و$\Sat{M}{!T(M,w)}$ بالبناء؛ فيلزم $\Sat{M}{!E(M,w)}$.
لكن $\Sat{M}{!E(M,w)}$ تكافئ وجود $n\in\Domain{M}=\Nat$
تكون فيه $M$ عند الدخل~$w$ في تهيئة توقف بعد~$n$ خطوات.
\end{explain}

\begin{defn}
نعرّف $!C(M,w,n)$ بأنها ال!!{sentence}
\[
\Obj Q_q(\num{m}, \num{n}) \land \Obj S_{\sigma_0}(\num{0}, \num{n})
\land \dots \land \Obj S_{\sigma_k}(\num{k}, \num{n}) \land
\lforall[x][(\num{k} < x \lif \Obj S_\TMblank(x, \num{n}))]
\]
وفيها $q$ حالة $M$ في الزمن~$n$، وتكون $M$ على الخانة~$m$
في الزمن~$n$؛ وتحمل الخانة~$i$ الرمز~$\sigma_i$ في الزمن~$n$
لكل $0\le i\le k$. وأما $k$ فهو الأكبر من رقمين: رقم آخر خانة
غير فارغة إلى اليمين في الزمن~$0$، ورقم أقصى خانة يمينًا زارها
الرأس خلال $n$ خطوة.
\end{defn}

\begin{lem}\ollabel{lem:halt-config-implies-halt}
متى بلغت $M$ عند الدخل~$w$ تهيئة توقف بعد $n$ خطوة، لزم
$!C(M,w,n)\Entails !E(M,w)$.
\end{lem}

\begin{proof}
لتتوقف $M$ عند الدخل~$w$ بعد $n$ خطوة. فثمة حالة~$q$ وخانة
$m$ ورمز~$\sigma$ تحقق ما يأتي:
\begin{enumerate}
\item تكون $M$ بعد $n$ خطوة في الحالة~$q$ على الخانة~$m$
  التي تحمل~$\sigma$.
\item لا تكون دالة الانتقال $\delta(q,\sigma)$ معرّفة.
\end{enumerate}
ووصف التهيئة $!C(M,w,n)$ يتضمن
$\Obj Q_q(\num{m},\num{n})$ و$\Obj S_\sigma(\num{m},\num{n})$.
ومن اجتماعهما تلزم $!E(M,w)$، أي:
\[
\lexists[x][\lexists[y][(\bigvee_{\tuple{q, \sigma} \in
      X}(\Obj Q_q(x, y) \land \Obj S_\sigma(x, y)))]]
\]
ووجه اللزوم أن $\Obj Q_q(\num{m},\num{n})\land
\Obj S_\sigma(\num{m},\num{n})$ تقتضي
$\bigvee_{\tuple{q,\sigma}\in X}(\Obj Q_q(\num{m},\num{n})\land
\Obj S_\sigma(\num{m},\num{n}))$، وذلك لأن
$\tuple{q,\sigma}\in X$.
\end{proof}

\begin{explain}
فمن توقف $M$ عند الدخل $w$ يلزم وجود~$n$ بحيث
$!C(M,w,n)\Entails !E(M,w)$. وننتقل إلى إثبات أنه لكل زمن~$n$
إلى أول تهيئة توقف يتحقق $!T(M,w)\Entails !C(M,w,n)$.
\end{explain}

\begin{lem}
\ollabel{lem:config}
لكل $n$ لم تتوقف فيه $M$ قبل تمام $n$ خطوة، يكون
$!T(M,w)\Entails !C(M,w,n)$.
\end{lem}

\begin{proof}
أما أساس الاستقراء، ففي الحالة $n=0$ نجد كل مقترن من
$!C(M,w,0)$ مقترنًا في $!T(M,w)$ نفسها؛ فالاستلزام حاصل.

وأما فرض الاستقراء فهو: إن لم تتوقف $M$ قبل الخطوة رقم $n$،
كان $!T(M,w)\Entails !C(M,w,n)$. والمطلوب أن نبين، ما دامت
$!C(M,w,n)$ ليست وصف توقف، أن $!T(M,w)\Entails !C(M,w,n+1)$.

لتكن $M$ بعد $n$ خطوة من ابتدائها عند $w$ في الحالة~$q$
على الخانة~$m$. فإن $M$، ما دامت لم تتوقف بعد~$n$ خطوة،
لا بد أن يحتوي برنامج~$M$ تعليمة من إحدى الصور الثلاث:

\begin{enumerate}
\item \ollabel{right} $\delta(q, \sigma) = \tuple{q', \sigma', \TMright}$

\item \ollabel{left} $\delta(q, \sigma) = \tuple{q', \sigma', \TMleft}$

\item \ollabel{stay} $\delta(q, \sigma) = \tuple{q', \sigma', \TMstay}$
\end{enumerate}

ولكل واحدة من هذه الصور حجة نوردها على الترتيب.

\begin{enumerate}
\item أما التعليمة من الصورة~\olref{right}، فبحسب
  \olref[rep]{defn:tm-descr}\olref[rep]{rep-right} تكون
\begin{align*}
& \lforall[x][\lforall[y][((\Obj Q_{q}(x,
  y) \land \Obj S_\sigma(x, y)) \lif {}]]\\
& \qquad (\Obj Q_{q'}(x',
  y') \land \Obj S_{\sigma'}(x, y') \land !A(x, y)))
\intertext{مقترنة في $!T(M,w)$. وتستلزم ال!!{sentence} الآتية
  (بالتخصيص الكلي، مع $\num{m}$ بدل~$x$ و$\num{n}$ بدل~$y$):}
& (\Obj Q_{q}(\num{m}, \num{n}) \land \Obj S_{\sigma}(\num{m},
\num{n})) \lif {}\\
& \qquad (\Obj Q_{q'}(\num{m}', \num{n}') \land
\Obj S_{\sigma'}(\num{m}, \num{n}') \land !A(\num{m}, \num{n})).
\intertext{ويفيد فرض الاستقراء أن $!T(M,w)\Entails !C(M,w,n)$؛ أي}
& \Obj Q_q(\num{m}, \num{n}) \land \Obj S_{\sigma_0}(\num{0}, \num{n})
\land \dots \land \Obj S_{\sigma_k}(\num{k}, \num{n}) \land \\
& \qquad\lforall[x][(\num{k} < x \lif \Obj S_\TMblank(x, \num{n}))]\\
\intertext{ولما كانت الخانة~$m$ بعد $n$ خطوة تحتوي~$\sigma$، كان
  المقترن المقابل $\Obj S_\sigma(\num{m},\num{n})$، ولذلك تستلزم هذه}
& \Obj Q_{q}(\num{m}, \num{n}) \land \Obj S_{\sigma}(\num{m},
   \num{n}).
\intertext{فنحصل الآن على}
& \Obj Q_{q'}(\num{m}', \num{n}') \land \Obj S_{\sigma'}(\num{m},
  \num{n}') \land {}\\
& \qquad\displaystyle\bigwedge_{\substack{0\leq i\leq k\\i\neq m}}
  \Obj S_{\sigma_i}(\num{i}, \num{n}') \land {}\\
& \qquad \lforall[x][(\num{k} < x \lif \Obj S_\TMblank(x, \num{n}'))]
\end{align*}
% تصحيح مصدر (C271-01): قُيّد حفظ رموز الشريط بالخانات i≠m؛ فالخانة m تُكتب فيها sigma'.
وبيان ذلك أن السطر الأول مأخوذ بقياس الفصل من تالي الشرط
المتقدم. وأما مقترنات السطر الأوسط، مع استثناء
$S_{\sigma_m}(\num{m},\num{n}')$، فيؤخذ كل واحد منها من
نظيره في $!C(M,w,n)$ مع $!A(\num{m},\num{n})$.

فإن كان $m<k$، كان $!T(M,w)\Proves\num{m}<\num{k}$
(\olref[rep]{prop:mlessk})، ويعطي تعدي~$<$
$\lforall[x][(\num{k}<x\lif\num{m}<x)]$. وإن كان $m=k$،
ثبت $\lforall[x][(\num{k}<x\lif\num{m}<x)]$ بالمنطق وحده.
فيؤخذ السطر الأخير من المقترن الموافق في $!C(M,w,n)$، ومن
$\lforall[x][(\num{k}<x\lif\num{m}<x)]$، ومن
$!A(\num{m},\num{n})$. وفي الحالة $m<k$ يكون الحاصل هو
$!C(M,w,n+1)$ بلا زيادة.

وأما إذا كان $m=k$، فقد بلغ الرأس عند تمام $n+1$ خطوة
الخانة~$k+1$، فصارت أقصى ما زاره يمينًا. ولهذا يزاد في
$!C(M,w,n+1)$ المقترن $\Obj S_\TMblank(\num{k}',\num{n}')$،
ويكون آخر مقترناتها
$\lforall[x][(\num{k}'<x\lif\Obj S_\TMblank(x,\num{n}'))]$.
فيلزمنا إثبات استلزام هاتين ال!!{sentence}s أيضًا.

قد حصل لنا $\lforall[x][(\num{k}<x\lif\Obj S_\TMblank(x,
  \num{n}'))]$؛ ومن تخصيصها
$\num{k}<\num{k}'\lif\Obj S_\TMblank(\num{k}',\num{n}')$.
والبديهية $\lforall[x][x<x']$ تعطي $\num{k}<\num{k}'$،
فيتم بقياس الفصل $\Obj S_\TMblank(\num{k}',\num{n}')$.

ثم إن $!T(M,w)\Proves\num{k}<\num{k}'$، فبضم بديهية تعدي
$<$ إلى ما تقدم نحصل على
$\lforall[x][(\num{k}'<x\lif\Obj S_\TMblank(x,\num{n}'))]$.
وتفصيل هذا الاستنتاج متروك للتمرين.

\item وأما التعليمة من الصورة~\olref{left}، فبحسب
  \olref[rep]{defn:tm-descr}\olref[rep]{rep-left}،
\begin{align*}
& \lforall[x][\lforall[y][((\Obj Q_{q}(x', y) \land \Obj
    S_{\sigma}(x', y)) \lif {}]]\\
& \qquad (\Obj Q_{q'}(x, y') \land \Obj
  S_{\sigma'}(x', y') \land !A(x', y))) \land {}\\
& \lforall[y][((\Obj Q_{q}(\num{0}, y) \land \Obj S_{\sigma}(\num{0},
    y)) \lif {}]\\
& \qquad (\Obj Q_{q'}(\num{0}, y') \land \Obj S_{\sigma'}(\num{0}, y')
  \land !A(\num{0}, y)))
\intertext{مقترنة في $!T(M,w)$. فإذا كان $m>0$، فليكن $l=m-1$
  (أي $m=l+1$). ويستلزم المقترن الأول من ال!!{sentence} السابقة ما يأتي:}
& (\Obj Q_{q}(\num{l}', \num{n}) \land \Obj S_{\sigma}(\num{l}', \num{n}))
\lif {} \\
& \qquad (\Obj Q_{q'}(\num{l}, \num{n}') \land \Obj S_{\sigma'}(\num{l}',
\num{n}') \land !A(\num{l}', \num{n})).
\intertext{وإلا فليكن $l=m=0$، وانظر إلى ال!!{sentence} الآتية التي
  يستلزمها المقترن الثاني:}
& ((\Obj Q_{q}(\num{0}, \num{n}) \land \Obj S_{\sigma}(\num{0}, \num{n})) \lif {}\\
& \qquad   (\Obj Q_{q'}(\num{0}, \num{n}') \land \Obj S_{\sigma'}(\num{0}, \num{n}') \land
!A(\num{0}, \num{n}))).
\intertext{وتستلزم كل من العبارتين}
&  \Obj Q_{q'}(\num{l}, \num{n}') \land \Obj S_{\sigma'}(\num{m},
  \num{n}') \land {}\\
&\qquad \displaystyle\bigwedge_{\substack{0\leq i\leq k\\i\neq m}}
  \Obj S_{\sigma_i}(\num{i}, \num{n}') \land {} \\
& \qquad  \lforall[x][(\num{k} < x
  \lif \Obj S_\TMblank(x, \num{n}'))]
\end{align*}
% تصحيح مصدر (C271-02): رُبط فرع أقصى اليسار بالحالتين q وq' في التعليمة محل البحث.
على الوجه المتقدم. ففي الحالة الأولى
$\num{l}'\ident\num{l+1}\ident\num{m}$، وفي الثانية
$\num{l}\ident\num{0}$. والحاصل هو $!C(M,w,n+1)$ نفسه.

\item وأما الحالة~\olref{stay} فبرهانها تمرين.
\end{enumerate}
فقد ثبت لكل~$n$ إلى أول تهيئة توقف أن
$!T(M,w)\Entails !C(M,w,n)$.
\end{proof}

\begin{prob}
تمم حجة الحالة~\olref[tur][und][ver]{stay} من برهان
\olref[tur][und][ver]{lem:config}.
\end{prob}

\begin{prob}
أنشئ !!a{derivation} نتيجته $\Obj S_{\sigma_i}(\num{i},\num{n}')$، ومقدماته
$\Obj S_{\sigma_i}(\num{i},\num{n})$ و$!A(m,n)$ (على فرض
$i\neq m$؛ أي إما $i<m$ وإما $m<i$).
\end{prob}

\begin{prob}
أنشئ !!a{derivation} نتيجته
$\lforall[x][(\num{k}'<x\lif\Obj S_\TMblank(x,\num{n}'))]$، ومقدماته
$\lforall[x][(\num{k}<x\lif\Obj S_\TMblank(x,\num{n}'))]$، و
$\lforall[x][x<x']$، و
$\lforall[x][\lforall[y][\lforall[z][((x<y\land y<z)\lif x<z)]]]$.)
\end{prob}

\begin{lem}
\ollabel{lem:valid-if-halt}
توقف $M$ عند الدخل~$w$ يقتضي أن تكون $!T(M,w)\lif !E(M,w)$ صادقة
منطقيًا.
\end{lem}

\begin{proof}
بمقتضى \olref{lem:config}، لكل زمن~$n$ إلى أول تهيئة توقف،
يكون الوصف~$!C(M,w,n)$ لتهيئة $M$ في الزمن~$n$ لازمًا
من~$!T(M,w)$. ولتتوقف $M$ بعد $k$ خطوات، ورأسها على
الخانة~$m$ لبعض $m\in\Nat$. عندئذ تصف $!C(M,w,k)$ تهيئة
توقف~$M$؛ فهي تتضمن المقترنين
$\Obj Q_q(\num{m},\num{k})$ و$\Obj S_\sigma(\num{m},\num{k})$،
مع كون $\delta(q,\sigma)$ غير معرّفة. فتعطي
\olref{lem:halt-config-implies-halt}
$!C(M,w,k)\Entails !E(M,w)$. وقد ثبت
$!T(M,w)\Entails !C(M,w,k)$، فينتج
$!T(M,w)\Entails !E(M,w)$، وتكون
$!T(M,w)\lif !E(M,w)$ صادقة منطقيًا.
\end{proof}

\begin{explain}
وبقي لإتمام التحقق عكس هذه النتيجة: من الصدق المنطقي لـ
$!T(M,w)\lif !E(M,w)$ يلزم أن $M$ تتوقف فعلًا عند
الدخل~$w$.
\end{explain}

\begin{lem}
\ollabel{lem:halt-if-valid}
من $\Entails !T(M,w)\lif !E(M,w)$ يلزم توقف $M$ عند
الدخل~$w$.
\end{lem}

\begin{proof}
لنأخذ ال!!{structure}~$\Struct M$ للغة $\Lang L_M$ ذات
المجال $\Nat$؛ نفسر فيها $\Obj 0$ بالعدد~$0$،
و$\prime$ بدالة الخلف، و$<$ بعلاقة الأصغر من. ونجعل تفسير
المحمولين $\Obj Q_q$ و~$\Obj S_\sigma$ على الوجه الآتي:
\begin{align*}
  \Assign{\Obj Q_q}{M} & =
\Setabs{\tuple{m, n}}{\begin{array}{ll}\text{بعد بدء $M$ عند~$w$ ومرور~$n$ خطوة،}\\ \text{تكون في الحالة $q$
  وتمسح الخانة~$m$}\end{array}} \\
\Assign{\Obj S_\sigma}{M} & = \Setabs{\tuple{m, n}}{\begin{array}{ll}
\text{بعد بدء $M$ عند~$w$ ومرور $n$ خطوة،}\\ \text{تحتوي الخانة~$m$ من شريط $M$
  الرمز~$\sigma$}\end{array}}
\end{align*}
فقد بنينا ال!!{structure}~$\Struct{M}$ لتطابق، خطوة بعد خطوة،
ما تفعله $M$ حقًا عند الدخل~$w$. ومن البناء يتضح
$\Sat{M}{!T(M,w)}$. فإن كان $\Entails !T(M,w)\lif !E(M,w)$،
لزم $\Sat{M}{!E(M,w)}$، أي:
\[
\Sat{M}{\lexists[x][\lexists[y][(\bigvee_{\tuple{q, \sigma} \in
      X}(\Obj Q_q(x, y) \land \Obj S_\sigma(x, y)))]]}.
\]
وحيث إن $\Domain{M}=\Nat$، فلا بد من $m,n\in\Nat$ يتحقق لهما
$\Sat{M}{\Obj Q_q(\num{m},\num{n})\land
\Obj S_\sigma(\num{m},\num{n})}$ لبعض~$q$ و~$\sigma$،
وتكون $\delta(q,\sigma)$ غير معرّفة. وهذا، بتعريف~$\Struct M$،
معناه أن $M$ تبدأ عند الدخل~$w$، ثم تكون بعد~$n$ خطوة في
الحالة~$q$ قارئةً الرمز~$\sigma$، ولا انتقال لها؛ أي إن~$M$
قد توقفت.
\end{proof}

\end{document}
